\begin{statementbox}
	\textbf{\textcolor{CStatement}{条件}}
	\begin{description}[style=nextline, leftmargin=2.8em, labelsep=0.5em, itemsep=0.35em]
		\item[\textbf{\textcolor{CStatement}{条件 1（函数形式）：}}] 设函数解析式为 $y = kx + b$（$k,b$ 为常数），定义域为 $\mathbb{R}$。
		\item[\textbf{\textcolor{CStatement}{条件 2（两点横坐标不同）：}}] 若已知直线上两点 $(x_1,y_1),(x_2,y_2)$ 求斜率，必须有 $x_1 \neq x_2$。
	\end{description}

	\textbf{\textcolor{CStatement}{结论}}
	\begin{description}[style=nextline, leftmargin=2.8em, labelsep=0.5em, itemsep=0.45em]
		\item[\textbf{\textcolor{CStatement}{结论一（斜率与截距）：}}]
			\textbf{\textcolor{CStatement}{直观描述：}}
			斜率由两点坐标差确定，纵截距是 $y$ 轴上的交点。
			\[
				k = \frac{y_2 - y_1}{x_2 - x_1} = \frac{\Delta y}{\Delta x},\quad b = y(0).
			\]

		\item[\textbf{\textcolor{CStatement}{结论二（单调性分类）：}}]
			\textbf{\textcolor{CStatement}{直观描述：}}
			斜率正负决定函数的增减。
			\[
				k > 0 \Rightarrow \text{递增},\quad k < 0 \Rightarrow \text{递减},\quad k = 0 \Rightarrow \text{常函数}.
			\]

		\item[\textbf{\textcolor{CStatement}{推广结论（斜率不存在）：}}]
			\textbf{\textcolor{CStatement}{直观描述：}}
			当直线垂直 $x$ 轴时，斜率不存在，无法使用两点式。
			\[
				x_1 = x_2 \;\Longrightarrow\; \text{斜率不存在},\; \text{直线方程为 } x = x_1.
			\]
	\end{description}
\end{statementbox}
